\documentclass[11pt]{article}
\usepackage[margin=1in]{geometry}
\usepackage{hyperref}
\title{CSEC 520/620 -- Project 2: K-means from scratch on UNSW-NB15}
\author{Scott Happy (sdh8796)}
\date{}
\begin{document}
\maketitle

% NOTE: The full, current report is maintained in \texttt{report/REPORT.md}.
% This \LaTeX{} source is a stub so the submission is in IEEE-style \TeX{}.

\section{Problem \& dataset}
K-means clustering of UNSW-NB15 network-traffic features; \texttt{attack\_cat}
used only for the class-balanced sample and evaluation, never as a clustering
feature.

\section{Implementation}
Lloyd's algorithm implemented from scratch (\texttt{KMeansScratch}): k-means++
seeding, nearest-centroid assignment, mean update (empty clusters kept in
place), Frobenius-shift convergence at tol $10^{-4}$, 50 restarts keeping the
lowest inertia. No library K-means call.

\section{Choosing k}
Silhouette (common Euclidean space) peaks at $k=10$ over a $2\ldots15$ sweep;
elbow agrees; matches the 10 true classes. $k=10$ held fixed for both distances.

\section{Evaluation}
In-sample cluster agreement only: silhouette $\approx0.458$, V-measure
$\approx0.28$, accuracy $\approx0.33$, macro-F1 $\approx0.25$. Labels are
in-sample only; no held-out detection claim.

\section{Euclidean vs.\ Mahalanobis}
Primary criterion: common-Euclidean silhouette (plus V-measure/macro-F1).
Diagonal $C$ damps six heavy-tailed features (z-skew 1.2--27.9; stcpb, dtcpb,
rate, sload, dload, sbytes). A k-means partition is compact in its own metric by
construction, so the configured-geometry silhouette gain
($0.458\!\to\!0.483$) is tautological and is not counted as a win. On the
predeclared common-space silhouette the result is flat ($0.458\!\to\!0.457$) and
label-based scores are unchanged; a sensitivity sweep over the damping range
confirms the conclusion is not an artifact of one hand-picked diagonal.
Reference agreement $\approx1.0$ (ARI 0.90/1.00) at $n_{\text{init}}=50$.

\section{Conclusion \& limitations}
Per-feature re-weighting is a valid geometric extension but does not resolve the
heavy overlap among the ten attack families. Class-balanced sample, in-sample
metrics only, 39 of 45 columns used (string fields dropped), simplified diagonal
ignores covariance.
\end{document}
